% Compile twice: lualatex 20001253.tex
\documentclass[11pt]{article}
\usepackage[margin=25mm]{geometry}
\usepackage{fontspec,unicode-math}
\setmainfont{Latin Modern Roman}
\setmathfont{Latin Modern Math}
\usepackage{amsmath,amssymb,mathtools}
\usepackage{xurl,hyperref}
\hypersetup{colorlinks=true,urlcolor=blue,linkcolor=blue}
\setcounter{secnumdepth}{0}
\setlength{\parindent}{0pt}
\setlength{\parskip}{5pt}
\setlength{\emergencystretch}{3em}
\newcommand{\sref}[2]{\hyperlink{source:#1:#2}{[S#2]}}
\newcommand{\cataloguescope}[1]{\paragraph{Source question.}#1}
\newcommand{\floor}[1]{\left\lfloor#1\right\rfloor}
\begin{document}
% UnsolvedMath ID 20001253; source catalogue code AIM-COMPUTATION-0091
\setcounter{section}{20001252}
\section{Perfect-number recognition and simultaneous polynomial divisors}\label{20001253}
{\small\sffamily UnsolvedMath ID 20001253 · Computer Science}
\subsection{Definitions and mathematical statement}

\paragraph{Shared notation.}Standard number systems are denoted by $\mathbb N,\mathbb Z,\mathbb Q,\mathbb R,\mathbb C$; index ranges, computational models and specialized notation follow the cited problem source and the definitions below.


Polynomial time means a bound polynomial in the finite input length; for an integer input, its binary length is proportional to its logarithm. A positive integer $N$ is perfect if $\sum_{d\mid N,\,d<N}d=N$, equivalently $\sigma(N)=2N$. The second question asks for integer arguments at which two given monic integer polynomials divide prescribed positive integers. Coprimality is over $\mathbb Q[X]$. Both recognition and bounded enumeration targets are retained.
\paragraph{Statement recorded in the supplied source.}
Question 3. Find a deterministic polynomial time algorithm for recognizing perfect numbers. (Lenstra) Given two monic polynomials f, g ∈ Z[X] with no common irreducible factor, and two integers n, m ∈ Z>0, find in polynomial time all x ∈ Z such that f (x) | m and g(x) | n, with

|x| < H (f, g, m, n ) for some function H. \sref{20001253}{2}
\paragraph{Scope qualification.}
The bound $H$ is not fixed by the source. A trivial tiny search bound would not settle the intended enumeration task; output-length and bound conventions need to be specified in any proposed algorithm.
\paragraph{Source status.}
The supplied catalogue labels this record partially solved. Its source-status metadata was checked against the supplied snapshot on 2026-10-11. This import does not independently establish current open status or mathematical correctness.
\subsection{Short English statement}
Can perfect integers be recognized deterministically in polynomial time, and can simultaneous polynomial-divisibility conditions be solved efficiently?
\subsection{Sources}
\begin{itemize}
\item[\hypertarget{source:20001253:1}{S1}] Ulam AI and the UnsolvedMath Contributors, UnsolvedMath: A Curated Collection of Open Mathematics Problems, record 20001253 (AIM-COMPUTATION-0091). Catalogue curation and metadata: CC BY 4.0. The supplied source question is retained; definitions and the English summary are editorial additions.
\url{https://huggingface.co/datasets/ulamai/UnsolvedMath}.
\item[\hypertarget{source:20001253:2}{S2}] AIM Workshop Problem Lists, original problem source linked by catalogue record AIM-COMPUTATION-0091
\url{https://aimath.org/WWN/primesinp/primesinp.pdf}.
\end{itemize}
\paragraph{Preparation and review.}
Original question attributed to the cited source and UnsolvedMath contributors. Codex (GPT-6; exact serving revision not exposed) assisted with definitions, English summaries, source inspection and formatting. No new human mathematical authorship or independent proof verification is asserted. Canonical ID 20001253 and source code AIM-COMPUTATION-0091 checked against the complete supplied registry and existing repository statements; known flagged repetitions and obvious identical conjectures excluded. Full mathematical-equivalence/literature review has not been performed. Statement preparation and final review followed contribution prompts C and D; automated checks verify structure and identity, not open status.
\label{end:20001253}
\end{document}
